\documentclass[11pt,a4paper]{article}
\usepackage[top=1.7cm,bottom=1.5cm,left=1.5cm,right=1.5cm,headsep=8pt,footskip=16pt]{geometry}
\usepackage{enumitem}
\usepackage{amssymb}
\usepackage{array}
\usepackage[table]{xcolor}
\usepackage{fancyhdr}
\usepackage{lastpage}
\usepackage{multicol}
\usepackage{polyglossia}
\setdefaultlanguage[calendar=gregorian, numerals=maghrib]{arabic}
\setotherlanguage{english}
\IfFontExistsTF{Amiri}{%
  \newfontfamily\arabicfont[Script=Arabic, Scale=1.08]{Amiri}
}{%
  \newfontfamily\arabicfont[Script=Arabic, Scale=1.08, Path=../fonts/, Extension=.ttf,
    UprightFont=*-Regular, BoldFont=*-Bold, ItalicFont=*-Italic,
    BoldItalicFont=*-BoldItalic]{Amiri}
}
\IfFontExistsTF{AlBattar}{%
  \newfontfamily\khatt[Script=Arabic, Scale=1.15]{AlBattar}
}{%
  \newfontfamily\khatt[Script=Arabic, Scale=1.15, Path=../fonts/,
    BoldFont=ae_AlBattar.ttf]{ae_AlBattar.ttf}
}
\renewcommand{\headrulewidth}{0.8pt}
\renewcommand{\footrulewidth}{0.8pt}
\fancyhf{}
\rhead{\khatt المدرسة العليا للأساتذة ببوزريعة}
\lhead{\khatt ملحقة بومرداس}
\lfoot{\khatt الأستاذ: د. بن عبد الرزاق يوسف}
\cfoot{\thepage/\pageref{LastPage}}
\pagestyle{fancy}
\parindent 0pt
\setlength{\columnsep}{0.6cm}
\setlength{\columnseprule}{0.4pt}
\setlist[enumerate,1]{label=\textbf{\arabic*.}, leftmargin=*, itemsep=5pt, topsep=2pt}
\begin{document}

\linespread{1.08}\selectfont
\begin{center}
{\LARGE\khatt\textbf{الاختبار الثاني عشر في الذكاء الاصطناعي}}\\[4pt]
{\normalsize وضعيات إشكالية ــ أسئلة متعددة الاختيارات ــ الفصول من 1 إلى 6}
\end{center}
\vspace{-2pt}\rule{\textwidth}{0.6pt}\\[2pt]
\textbf{التعليمة:} يتضمن الاختبار خمس وضعيات إشكالية. اقرأ كل وضعية بعناية، ثم اختر التحليل أو التصرف الأنسب
من بين الاختيارات الخمسة؛ لكل وضعية إجابة صحيحة واحدة فقط.
اكتب \textbf{حرف} الإجابة الصحيحة في جدول الإجابات في آخر الاختبار.
لكل وضعية 4 نقاط (المجموع 20 نقطة).
\vspace{4pt}

\noindent
{\large\khatt\textbf{الوضعية 1}}\hfill{\small(الفصل 2)}\par\nopagebreak\vspace{3pt}
\setlength{\fboxsep}{7pt}\colorbox{gray!10}{\parbox{\dimexpr\linewidth-14pt}{كلّف أستاذ ثلاثة طلبة بتصنيف ألف تعليق لتلاميذ على نصوص أدبية إلى «إيجابي» و«سلبي» و«محايد»، لتدريب نموذج. تبيّن لاحقًا أن الطلبة اختلفوا في تصنيف ثلث التعليقات، وأن كل واحد اعتمد فهمه الخاص للمحايد. ويرغب في بدء التدريب في أقرب وقت. ما أثر ذلك على النموذج، وما الحل؟}}\par\nopagebreak
\begin{itemize}[leftmargin=1.0cm, labelsep=6pt, itemsep=3pt, topsep=4pt, parsep=0pt]
  \item[\textbf{أ)}] تضارب التسميات يربك النموذج فيتعلم حدودًا غامضة بين الأصناف؛ والحل وضع دليل تصنيف واضح بأمثلة، وتدريب المصنّفين عليه، ومراجعة التعليقات المختلف فيها قبل التدريب.
  \item[\textbf{ب)}] لا أثر لذلك، لأن النموذج يكتشف التصنيف الصحيح بنفسه مهما تضاربت التسميات، فالكمية الكبيرة من التعليقات تعوّض أي اختلاف بين المصنّفين في فهم الأصناف.
  \item[\textbf{ج)}] الحل أن يعتمد تصنيف الطالب الأكثر إنتاجًا، لأن الأسرع في التصنيف هو الأكثر فهمًا للأصناف، دون حاجة إلى دليل مشترك أو مراجعة التعليقات المختلف فيها.
  \item[\textbf{د)}] التضارب يربك النموذج، والحل مضاعفة عدد التعليقات إلى ألفين بالطريقة نفسها، لأن زيادة البيانات تلغي أثر الاختلاف، دون حاجة إلى توحيد فهم المصنّفين.
  \item[\textbf{هـ)}] الحل حذف صنف «محايد» نهائيًا، لأن وجود ثلاثة أصناف يربك كل النماذج، والنماذج لا تستطيع التعلم إلا من صنفين متقابلين مهما كانت جودة التسميات.
\end{itemize}
\par\vspace{6pt}
\noindent
{\large\khatt\textbf{الوضعية 2}}\hfill{\small(الفصل 3)}\par\nopagebreak\vspace{3pt}
\setlength{\fboxsep}{7pt}\colorbox{gray!10}{\parbox{\dimexpr\linewidth-14pt}{تريد أستاذة إنتاج تمارين نحوية بثلاثة مستويات لقسم متفاوت: داعم ومتوسط ومتقدم. طلبت «ثلاثة تمارين مختلفة الصعوبة عن المفعول به»، فجاءت التمارين متقاربة جدًا، والفرق بينها في عدد الجمل فقط. وتريد فروقًا حقيقية في المهارة المطلوبة. ولا تملك وقتًا كافيًا لكتابة التمارين بنفسها كل أسبوع، وتريد طلبًا واحدًا تعيد استعماله. أي طلب أنسب؟}}\par\nopagebreak
\begin{itemize}[leftmargin=1.0cm, labelsep=6pt, itemsep=3pt, topsep=4pt, parsep=0pt]
  \item[\textbf{أ)}] تحدد المهارة لكل مستوى، لكنها تطلب أن تكون الجمل نفسها في المستويات الثلاثة، لأن توحيد الجمل يمنع الأداة من تغيير الصعوبة، فيبقى الفرق في عدد الأسئلة وحده.
  \item[\textbf{ب)}] تطلب تمرينًا واحدًا صعبًا وتوزعه على الجميع، لأن التمرين الصعب يرفع مستوى الضعاف ويناسب المتقدمين، والتمييز بين المستويات يضر بتكافؤ الفرص في القسم.
  \item[\textbf{ج)}] تحدد لكل مستوى المهارة المطلوبة: الداعم يتعرّف المفعول به في جمل بسيطة، والمتوسط يعرّبه ويميّز أنواعه، والمتقدم يوظّفه في فقرة؛ مع مثال لكل مستوى وتصحيح منفصل.
  \item[\textbf{د)}] تكتب «ثلاثة تمارين: سهل جدًا، ومتوسط جدًا، وصعب جدًا»، لأن تشديد الصفات يجعل الأداة توسّع الفروق بين المستويات تلقائيًا، دون تحديد المهارة المطلوبة في كل مستوى.
  \item[\textbf{هـ)}] تطلب ثلاثين تمرينًا وترتّبها حسب طولها، لأن طول التمرين هو المعيار الحقيقي لصعوبته، ثم تختار أقصرها للداعم وأطولها للمتقدم دون مراجعة المهارة المطلوبة.
\end{itemize}
\par\vspace{6pt}
\noindent
{\large\khatt\textbf{الوضعية 3}}\hfill{\small(الفصل 4)}\par\nopagebreak\vspace{3pt}
\setlength{\fboxsep}{7pt}\colorbox{gray!10}{\parbox{\dimexpr\linewidth-14pt}{طلب طالب من أستاذه رسالة توصية لبرنامج ماستر، فكتب الأستاذ مسودة بأداة ذكاء اصطناعي بعد أن زوّدها بمعلومات عامة عن الطالب، فجاءت الرسالة مليئة بمدائح عامة ومعلومات عن «مشاريع بحثية» لم ينجزها الطالب. والطالب ينتظر الرسالة قبل نهاية الأسبوع، وموعد الترشح يقترب. ما التصرف الأسلم أخلاقيًا؟}}\par\nopagebreak
\begin{itemize}[leftmargin=1.0cm, labelsep=6pt, itemsep=3pt, topsep=4pt, parsep=0pt]
  \item[\textbf{أ)}] يحذف المشاريع غير المنجزة ويُبقي المدائح العامة كما هي، لأن المدائح لا تحتاج إلى شواهد، والرسالة المليئة بالصفات الإيجابية أنفع للطالب من الرسالة القائمة على وقائع محددة.
  \item[\textbf{ب)}] يحذف كل معلومة غير صحيحة، ويستبدل المدائح العامة بشواهد محددة يعرفها شخصيًا عن الطالب، لأن الأستاذ يتحمل مسؤولية كل ما يوقّعه، ورسالة التوصية شهادة لا نص إنشائي.
  \item[\textbf{ج)}] يرسل الرسالة كما هي مع عبارة «كُتبت بمساعدة الذكاء الاصطناعي»، لأن الإفصاح يرفع عنه المسؤولية عن المعلومات غير الصحيحة، ويبقى التوقيع إجراءً شكليًا.
  \item[\textbf{د)}] يطلب من الطالب نفسه مراجعة الرسالة وتعديلها كما يشاء، ثم يوقّعها دون قراءة، لأن الطالب أعرف بإنجازاته، وتوقيع الأستاذ يكفي لإثبات صحة ما يكتبه الطالب عن نفسه.
  \item[\textbf{هـ)}] يرسل الرسالة كما هي، لأن المبالغة مقبولة في رسائل التوصية، والمعلومات الإضافية تقوّي ملف الطالب، والأداة هي المسؤولة عن أي معلومة غير دقيقة وردت في الرسالة، فهذا ما يفعله كثير من الموقّعين.
\end{itemize}
\par\vspace{6pt}
\noindent
{\large\khatt\textbf{الوضعية 4}}\hfill{\small(الفصل 5)}\par\nopagebreak\vspace{3pt}
\setlength{\fboxsep}{7pt}\colorbox{gray!10}{\parbox{\dimexpr\linewidth-14pt}{طلبت طالبة من أداة ذكاء اصطناعي «أهم المصادر الأولية عن ثورة التحرير الجزائرية»، فقدّمت قائمة تضم كتبًا تاريخية حديثة ومقالات صحفية ومذكرات لبعض المجاهدين، دون تمييز بينها. وأستاذها يشترط الاعتماد على مصادر أولية. ولم تتعلم بعد طريقة التمييز بين أنواع المصادر، وتعتمد على الأداة في ذلك. كيف تتعامل مع القائمة؟}}\par\nopagebreak
\begin{itemize}[leftmargin=1.0cm, labelsep=6pt, itemsep=3pt, topsep=4pt, parsep=0pt]
  \item[\textbf{أ)}] تعتمد المقالات الصحفية الحديثة وحدها، لأنها الأحدث، والمصدر الأولي هو الأحدث نشرًا دائمًا، ثم تتحقق من وجودها فقط دون النظر في تاريخ كتابتها أو صلتها بالأحداث.
  \item[\textbf{ب)}] تحذف المذكرات من القائمة، لأن مذكرات الأشخاص آراء ذاتية لا تُعدّ مصادر أولية، وتعتمد الكتب الحديثة وحدها، لأنها الأوثق والأقرب إلى شرط الأستاذ في البحث.
  \item[\textbf{ج)}] تطلب من الأداة تمييز الأولي من الثانوي، وتعتمد التصنيف الجديد دون أي تحقق، لأن الأداة تصحّح أخطاءها عند الطلب، ولا يمكن أن تختلق عناوين في المرة الثانية.
  \item[\textbf{د)}] تعتمد القائمة كما هي، لأن الأداة صنّفتها مصادر أولية، والأدوات تميّز المصادر الأولية من الثانوية بدقة، فلا حاجة إلى التحقق من نوع كل مصدر أو من وجوده، مهما كان نوعها أو تاريخها.
  \item[\textbf{هـ)}] تميّز بنفسها بين المصادر الأولية كالمذكرات والوثائق المعاصرة للأحداث، والثانوية كالدراسات الحديثة، وتتحقق من وجود كل عنوان وبياناته، لأن الأداة قد تخلط بين النوعين أو تختلق عناوين.
\end{itemize}
\par\vspace{6pt}
\noindent
{\large\khatt\textbf{الوضعية 5}}\hfill{\small(الفصل 6)}\par\nopagebreak\vspace{3pt}
\setlength{\fboxsep}{7pt}\colorbox{gray!10}{\parbox{\dimexpr\linewidth-14pt}{يريد أستاذ تحويل محاضراته المسجلة إلى نصوص مكتوبة بأداة تحويل الصوت إلى نص، ثم توزيعها على الطلبة. جاءت النصوص جيدة عمومًا، لكن فيها أخطاء في أسماء الأعلام والمصطلحات البلاغية، وبعض الآيات والأبيات منقولة بألفاظ قريبة لا مطابقة. ويريد الطلبة النصوص قبل الامتحان بأسبوع. ما الإجراء الأنسب قبل التوزيع؟}}\par\nopagebreak
\begin{itemize}[leftmargin=1.0cm, labelsep=6pt, itemsep=3pt, topsep=4pt, parsep=0pt]
  \item[\textbf{أ)}] يتخلى عن توزيع النصوص المكتوبة، لأن أخطاء الأداة تجعلها عديمة الفائدة، ويكتفي بالتسجيلات الصوتية، لأن النصوص المحوَّلة آليًا لا يمكن تصحيحها بجهد معقول.
  \item[\textbf{ب)}] يوزّع النصوص كما هي مع تنبيه عام إلى احتمال وجود أخطاء، لأن التنبيه يكفي ليحمي الطلبة من الأخطاء، ولا يلزم الأستاذ تصحيح الآيات والأبيات بنفسه.
  \item[\textbf{ج)}] يصحّح أسماء الأعلام فقط، لأنها الأخطاء الوحيدة المهمة، أما الآيات والأبيات فتكفي ألفاظها القريبة ما دام المعنى العام محفوظًا في النص المكتوب.
  \item[\textbf{د)}] يراجع النصوص مع التسجيلات، ويصحّح الأعلام والمصطلحات، ويقابل الآيات بالمصحف والأبيات بالدواوين حرفيًا، لأن الأداة تكتب الألفاظ المرجّحة صوتيًا ولغويًا، لا المطابقة بالضرورة.
  \item[\textbf{هـ)}] يعيد تحويل التسجيلات بأداة أخرى، ويوزّع النسخة التي تتفق فيها الأداتان، لأن اتفاق أداتين يضمن مطابقة الآيات والأبيات دون حاجة إلى مقابلتها بالمصادر.
\end{itemize}
\par\vspace{6pt}
\begin{center}
{\khatt\textbf{جدول الإجابات}}\\[3pt]
\renewcommand{\arraystretch}{1.6}
\begin{tabular}{|>{\centering\arraybackslash}m{1.9cm}|*{5}{>{\centering\arraybackslash}m{1.2cm}|}}
\hline
\textbf{رقم السؤال} & \textbf{1} & \textbf{2} & \textbf{3} & \textbf{4} & \textbf{5} \\ \hline
\textbf{الإجابة} &  &  &  &  &  \\[8pt] \hline
\end{tabular}

\vspace{8pt}
\renewcommand{\arraystretch}{2.2}
\begin{tabular}{|>{\centering\arraybackslash}m{7.5cm}|>{\centering\arraybackslash}m{3.2cm}|>{\centering\arraybackslash}m{3.2cm}|}
\hline
\textbf{الاسم واللقب} & \textbf{الفوج} & \textbf{العلامة} \\ \hline
 & & \hfill /20 \\ \hline
\end{tabular}
\end{center}
\end{document}
